% !TeX root = ../main.tex

\chapter{Architettura del Sistema e Osservabilità}
\label{cap:architettura_osservabilita}

L'architettura è basata su \texttt{Kind} (Kubernetes IN Docker), che avvia cluster Kubernetes locali in container Docker, replicando scheduler, networking e controller di un cluster reale in un ambiente facilmente ricreabile per validazione e test. Kubernetes orchestra i container organizzandoli in \texttt{Pod}, mantiene lo stato desiderato tramite \texttt{Deployment} e \texttt{ReplicaSet}, espone i servizi con \texttt{Service} e fornisce funzionalità native di autoscaling e self-healing. Docker fornisce la base di containerizzazione che confeziona applicazioni e dipendenze in immagini portabili, garantendo isolamento, avvii rapidi e coerenza tra ambienti.

\section{Design a Microservizi}
\label{sec:design_microservizi}

Il sistema separa responsabilità e domini di failure in tre blocchi indipendenti, ciascuno in Pod distinti con \texttt{Service} per la service discovery via DNS interno (\texttt{influxdb-service}, \texttt{mosquitto-service}):
\begin{itemize}
    \item \textbf{Data Layer} (InfluxDB): persistenza di metriche e serie temporali, stateful con PVC.
    \item \textbf{Message Broker} (Mosquitto/MQTT): canale asincrono che disaccoppia produttori e consumatori, permettendo riavvii e scaling del layer logico senza spezzare l'ingest.
    \item \textbf{Logic Layer} (\texttt{central\_server}, MCP, agenti): elaborazione e API di business.
\end{itemize}
Gli script Python sono eseguiti come processo principale del container tramite il \texttt{command} del Deployment; lo scheduler decide il nodo, il campo \texttt{replicas} distribuisce le istanze su nodi diversi mantenendo invariata l'interfaccia grazie ai \texttt{Service}.

\section{Simulazione del Mondo Fisico (Digital Twin)}
\label{sec:digital_twin}

I digital twin di droni e client sono generatori controllati di telemetria e domanda. Il twin del drone (\texttt{drone\_simulator.py}) mantiene stato interno (posizione, batteria, usura, stato operativo), pubblica telemetria su \texttt{telemetry/drones} e ascolta comandi su \texttt{comandi/<DRONE\_ID>}. Il twin del client (\texttt{client\_simulator.py}) genera carico business pubblicando ordini su \texttt{business/ordini/nuovi}.

La moltiplicazione della flotta è ottenuta scalando le repliche del Deployment \texttt{drone-simulator}: poiché \texttt{DRONE\_ID} è derivato dal \texttt{metadata.name} del Pod, ogni replica corrisponde a un drone digitale distinto. Lo scaling è invocabile anche dall'agente tramite il tool MCP \texttt{scale\_drone\_deployment} (patch su \texttt{deployment/scale}), con i guardrail descritti nel Capitolo~\ref{cap:safety_guardrails}.

\section{Strategia di Osservabilità}
\label{sec:osservabilita}

L'osservabilità si fonda su tre pilastri:
\begin{enumerate}
    \item \textbf{Dashboard web centralizzata} ospitata dal \texttt{central\_server} con endpoint REST Flask (\texttt{/api/status}, \texttt{/api/orders}, \texttt{/api/drones}, \texttt{/api/completed}); il client esegue polling \texttt{fetch()} ogni 5\,s.
    \item \textbf{Log strutturati e audit trail:} \texttt{audit\_actions.jsonl} (JSON Lines, traccia ogni azione di scrittura/rimediazione dell'IA: timestamp, tool, payload), \texttt{pending\_approvals.jsonl} (ciclo di vita delle richieste HITL con giustificazione LLM) e log di runtime con tracking dei token consumati per ciascun ciclo.
    \item \textbf{Serie temporali su InfluxDB:} bucket \texttt{iot\_data} con i measurement \texttt{drone\_telemetry} e \texttt{business\_orders}; gli LLM analizzano time-window di 5/60 minuti tramite i tool MCP \texttt{get\_drones\_telemetry} e \texttt{get\_pending\_orders}.
\end{enumerate}

Il registro \texttt{audit\_actions.jsonl}, in particolare, traccia per ogni operazione di rimediazione il timestamp, il tipo di tool invocato e l'esatto payload, fornendo l'\textit{auditability} richiesta dal dominio \textit{high-stakes}. In parallelo, l'orchestratore logga per ciascuna iterazione il consumo di token restituito dai modelli, metrica cruciale per la verifica dei tetti di spesa discussi nel Capitolo~\ref{cap:analisi_economica}.
